\chapter[Leave-one-generation-rule-out cross-validation]{Leave-one-generation-\\rule-out cross-validation}
\label{ch:loro}

\begin{saybox}[title=What to say at the start of this section]
Explain the procedure in plain words before any notation: ``we take away
one entire generation rule, train on the others, and ask the model to
predict the rule it has never seen --- then repeat for every rule.'' Then
state what this does and does not measure, because the panel will ask, and
volunteering it is better than conceding it.
\end{saybox}

\section{The procedure}

Let a specialist's rows be $\{(\mathbf{x}_i, y_i, g_i)\}_{i=1}^{n}$, where
$g_i \in \mathcal{G}$ is the row's \texttt{source\_rule\_id} and
$|\mathcal{G}| = G$. For each rule $g \in \mathcal{G}$:

\begin{enumerate}[leftmargin=1.5em, itemsep=1pt]
\item Define the held-out set $\mathcal{T}_g = \{ i : g_i = g \}$ and the
  training set $\mathcal{R}_g = \{ i : g_i \neq g \}$.
\item Fit the preprocessing transform on $\mathcal{R}_g$ only.
\item Fit the model on $\mathcal{R}_g$.
\item Predict every row in $\mathcal{T}_g$.
\item Record $\rsq$, MAE and RMSE on $\mathcal{T}_g$.
\end{enumerate}

With $n_{\text{splits}} = G$, \texttt{GroupKFold} grouped on
\texttt{source\_rule\_id} yields exactly leave-one-rule-out: each fold's
test set is precisely one rule. Verified in the implementation
(\evd{scripts/experiments/leave\_rule\_out\_row\_level.py}): the fold loop
asserts a single unique rule per test partition, the preprocessor is
constructed and fitted \emph{inside} the loop on the training rows only,
and hyperparameters are fixed a priori rather than tuned per fold --- so
no selection information crosses the boundary.

The complete evaluation comprises $56$ folds per model
($7+11+12+8+8+10$ rules across the six specialists), executed for Random
Forest, LightGBM and XGBoost, yielding $168$ fold results and
$102\,000$ row-level out-of-fold predictions.

\section{What this measures --- stated precisely}

\begin{limitbox}[title=Scope of the claim]
Leave-rule-out cross-validation measures \textbf{transfer across synthetic
generation rules}: the ability to predict rows produced by a generative
mechanism absent from training, using other synthetic mechanisms as the
training signal.

It does \textbf{not} measure transfer to independent experimental
observations. A high leave-rule-out $\rsq$ is a necessary condition for
real-world generalization, not a sufficient one. \Cref{ch:limits}
enumerates the dependencies that survive this protocol, and
\Cref{ch:external} presents the separate external evidence.
\end{limitbox}

\section{The metrics, and why their summaries disagree}

For a held-out set $\mathcal{T}$ with mean $\bar{y}_{\mathcal{T}}$,
\begin{align}
\rsq &= 1 - \frac{\sum_{i \in \mathcal{T}} (y_i - \hat{y}_i)^2}
                  {\sum_{i \in \mathcal{T}} (y_i - \bar{y}_{\mathcal{T}})^2}
       = 1 - \frac{\text{MSE}}{s_y^2}, \\
\text{MAE} &= \frac{1}{|\mathcal{T}|}\sum_{i \in \mathcal{T}} |y_i - \hat{y}_i|,
\qquad
\text{RMSE} = \sqrt{\frac{1}{|\mathcal{T}|}\sum_{i \in \mathcal{T}} (y_i - \hat{y}_i)^2}.
\end{align}

MAE and RMSE are in days and are directly interpretable: an MAE of $3$ days
means the average prediction misses by three days. $\rsq$ is a
\emph{relative} measure --- it compares the model against predicting the
held-out set's own mean. This distinction is the source of nearly every
counter-intuitive number in \Cref{ch:folds}.

\subsection{Four different summaries of the same folds}

\begin{description}[leftmargin=1.6em, style=nextline, font=\bfseries\color{dgreen}]
\item[Mean fold $\rsq$] $\frac{1}{G}\sum_g \rsq_g$. Weights every rule
  equally regardless of size, so an $8$-row rule counts as much as a
  $5\,654$-row rule. A single catastrophic small fold dominates it.
\item[Median fold $\rsq$] Robust to those outliers; describes the typical
  rule. Hides the failures entirely, which is its own danger.
\item[Pooled out-of-fold $\rsq$] Concatenate all out-of-fold predictions,
  then compute $\rsq$ once against the global mean. Weights by row count
  and is the closest analogue to the original headline number.
\item[MAE in days] Independent of any variance denominator and therefore
  immune to the small-fold pathology.
\end{description}

All four are reported throughout this document, because reporting only one
would misrepresent the result in one direction or the other.

\subsection{Why a small fold can produce a large negative $\rsq$}

If the held-out rule's own target variance $s_y^2$ is small, the
denominator is small, and a modest absolute error produces a large ratio.
Concretely, from the results in \Cref{ch:folds}: the fold
\texttt{HMMC\_NAT|NATAMYCIN\_REVIEW} has $n = 17$ rows, a target standard
deviation of $1.78$ days, a Random Forest MAE of $1.99$ days --- and
$\rsq = -0.335$. A two-day average error would be excellent in any
practical sense; it is negative in $\rsq$ terms only because the rule's
rows all lie within a few days of each other.

\begin{findingbox}[title=Reporting implication]
A negative fold $\rsq$ on a small, low-variance rule is not evidence of a
catastrophic practical error, and must not be presented as one. Equally, a
high median $\rsq$ is not evidence that no fold failed. Both summaries are
reported together throughout, alongside the error in days.
\end{findingbox}

\fig{fig22_mae_vs_r2}{0.9}{%
  Every Random Forest fold, plotted as MAE in days against fold $\rsq$.
  Marker size is fold row count; colour is the held-out rule's target
  standard deviation. The folds in the lower-left --- small error in days,
  poor $\rsq$ --- are the low-variance rules.}

\fig{fig23_r2_variance_decomposition}{0.9}{%
  The $\rsq$ ratio decomposed for the sixteen weakest folds: the target
  variance $s_y^2$ that forms the denominator against the mean squared
  error that forms the numerator. A fold goes negative exactly when the
  burgundy bar exceeds the green one.}

\section{Headline result across all seventeen models}

\tbl{tab_loro_full_zoo}{All seventeen evaluated configurations under
leave-one-rule-out cross-validation, averaged over the six specialists.}

\fig{fig13_zoo_leave_rule_out_ranking}{0.92}{%
  Leave-rule-out ranking. The three production tree ensembles separate
  clearly from every other family once the generation rule is withheld ---
  a separation that was invisible under the original protocol.}

\fig{fig14_holdout_vs_loro_slope}{0.82}{%
  The same models under the two protocols. Nothing about the models
  changed between the left and right columns; only the evaluation did.}

\begin{findingbox}[title=Effect of the protocol change]
Under the original protocol the seventeen models were bunched together,
with a linear model within two points of the production ensembles. Under
leave-rule-out they separate sharply:

\begin{itemize}[leftmargin=1.2em, itemsep=1pt]
\item LightGBM: mean $\rsq = 0.627$, median $0.780$
\item Random Forest: mean $\rsq = 0.582$, median $0.791$
\item XGBoost: mean $\rsq = 0.550$, median $0.719$
\item Ridge regression: mean $\rsq = -0.874$, median $-0.798$
\item $k$-NN ($k=5$): mean $\rsq = -3.412$, median $-1.537$
\item Dummy mean baseline: mean $\rsq = -15.637$
\end{itemize}

Ridge regression, which scored $0.910$ on soft/general under the original
protocol, is worse than useless when the rule is withheld. The capacity of
the tree ensembles turns out to matter a great deal --- but only under an
evaluation capable of detecting it.
\end{findingbox}

\fig{fig15_loro_heatmap_mean}{0.9}{%
  Leave-rule-out mean $\rsq$ for every model and specialist, clipped at
  $-2$ for legibility.}

\fig{fig16_loro_heatmap_median}{0.9}{%
  The same grid using median fold $\rsq$. Comparing with the previous
  figure isolates the specialists where the mean is distorted by a small
  number of extreme folds.}

\fig{fig17_mean_vs_median}{0.58}{%
  Mean against median fold $\rsq$ for every model--specialist combination.
  Points far below the diagonal indicate a distribution dominated by a few
  catastrophic folds.}

\section{Per-specialist results for the production ensembles}

\tbl{tab_loro_ensembles_by_specialist}{Leave-rule-out results for the three
production tree ensembles, broken out by specialist.}

\section{Pooled out-of-fold performance}

Because every row is predicted exactly once when it is held out, the
out-of-fold predictions can be concatenated and scored as a single set.
This weights by row count rather than by rule and is the most direct
counterpart to the original headline figure.

\tbl{tab_pooled_oof}{Pooled out-of-fold metrics: all folds concatenated
per specialist, then scored once. Computed from the $102\,000$ row-level
predictions.}

\fig{fig38_pooled_oof_r2}{0.95}{%
  Pooled out-of-fold $\rsq$ by specialist and model.}

\begin{findingbox}[title=The number to quote]
For the general-shelf-life task, pooled out-of-fold $\rsq$ under
leave-rule-out is $0.960$ (soft), $0.940$ (semi-hard) and $0.669$ (hard)
for LightGBM, against original context-holdout figures of $0.971$, $0.920$
and $0.932$.

Soft and semi-hard barely move; hard falls by more than twenty-five
points. That asymmetry is not noise, and \Cref{ch:folds} identifies its
specific cause --- a single held-out rule representing a spoilage
mechanism with a completely different time scale.
\end{findingbox}

\fig{fig37_error_ecdf}{0.78}{%
  Cumulative distribution of absolute out-of-fold error across all
  $34\,000$ rows per model. Median absolute error is under four days for
  all three ensembles; the distributions separate only in the upper tail.}
